\section{Introduction}
\label{sec:Introduction}

This manual describes the Raft library from two developer viewpoints.

The first viewpoint is the application developer who wants to add cluster
capability to a domain application. In that role the developer should not need
to understand every election rule or replication edge case. The important
questions are where application state enters the replicated log, how reads are
made safe, how commands and snapshots are encoded, and what the Java, C++ and
Rust integration surfaces expect.

The second viewpoint is the Raft machinery developer who maintains the library
itself. In that role the developer must reason about terms, voting, leader
state, log matching, snapshots, read leases, joint consensus, wire compatibility
and cross-language parity. The Java, C++ and Rust implementations are intended to be
functionally equivalent at the Raft boundary even though the language idioms,
runtime frameworks and internal structure differ.

The library exists to let applications run in a Raft cluster without embedding
domain knowledge in the consensus machinery. Reference-data distribution is an
important example: writes are relatively rare and governed, while lookups are
high-volume and need to happen close to consumers. The same library shape can
also support other master-data or configuration-data applications where a
central management path must feed fast replicated lookup state.

\subsection{What this manual covers}

The manual is organized around the boundary between domain application code and
Raft internals:

\begin{itemize}
\item Section~\ref{sec:GettingStarted} walks through building and running a cluster.
\item Section~\ref{sec:SystemOverview} summarizes the architecture and module layout.
\item Section~\ref{sec:ApplicationDeveloper} explains how an application uses the library.
\item Sections~\ref{sec:JavaLibrary},~\ref{sec:CppLibrary} and~\ref{sec:RustLibrary} describe the Java, C++ and Rust integration points.
\item Section~\ref{sec:RaftInternals} explains the consensus machinery for maintainers.
\item Section~\ref{sec:WireStorageSnapshots} covers protocol, persistence and snapshots.
\item Section~\ref{sec:TestingAndOperations} explains test strategy and operational checks.
\item Section~\ref{sec:JepsenTesting} documents the Jepsen fault-injection framework.
\end{itemize}

\subsection{Terminology}

This manual uses \emph{application} for the domain-specific code that owns
business state and command encoding. It uses \emph{Raft machinery} for elections,
log replication, membership, snapshot metadata, read-safety decisions and
transport handling. The boundary is deliberate: applications define what a
command means; Raft decides when that command is committed and safe to apply.

\section{Design intent}
\label{sec:DesignIntent}

The project should be read as a library project, not as a single reference-data
product. Reference data is a useful and important application, but the reusable
asset is the cluster capability:

\begin{itemize}
\item replicate a deterministic state machine,
\item preserve safety across leadership changes,
\item expose typed client command and query requests,
\item handle cluster membership through replicated configuration entries,
\item support Java, C++ and Rust applications through equivalent protocol behavior.
\end{itemize}

The Java implementation is the most comprehensive implementation. The C++ and
Rust implementations are developed as peer implementations and should be brought
up to the same behavioral maturity where differences affect Raft semantics.
Domain examples may differ temporarily, but the Raft machinery should not.
